\documentclass[11pt,a4paper]{article}

\usepackage[margin=1in]{geometry}
\usepackage{amsmath,amssymb}
\usepackage{graphicx}
\usepackage{booktabs}
\usepackage{siunitx}
\usepackage{hyperref}
\usepackage{caption}

\title{A Design-Oriented Flapping-Wing Aerodynamic Model\\(LaTeX rewrite with original equation numbering)}
\author{}
\date{}

\begin{document}
\maketitle

\begin{abstract}
This document rewrites the paper into a clean, copyable \LaTeX{} form. The prose is paraphrased for clarity, while the governing equations follow the original notation and equation numbers.
\end{abstract}

\tableofcontents

% =========================================================
\section{Overview}
We model the aerodynamic loads of harmonically flapping wings using a modified strip-theory approach.
Each spanwise strip is treated as a 2D section with finite-span unsteady corrections. The model predicts
instantaneous and cycle-averaged lift, thrust, input power, and propulsive efficiency, and includes
(i) circulatory and apparent-mass normal forces, (ii) chordwise contributions from camber, leading-edge suction
(with efficiency), and skin friction, and (iii) a localized post-stall (separated-flow) approximation with a
crossflow-drag normal-force model.

% =========================================================
\section{Method of Analysis}

\subsection{Section kinematics and basic normal force}
The section kinematics are defined using the leading edge as the reference point.
The motion of a section consists of a plunging velocity $\dot{h}$ and a pitch angle $\theta$.
Note that $\dot{h}$ is not necessarily perpendicular to the mean-stream velocity $U$; for root-flapping motion,
the displacement $h$ is perpendicular to the flapping axis. The wing aspect ratio is assumed sufficiently large
so that, in the mean-stream direction, the flow over each section may be treated as essentially chordwise.

The circulatory normal force on a strip is
\begin{equation}
dN_c=\frac{\rho U V}{2}\,C_n(y)\,c\,dy
\tag{1}\label{eq:1}
\end{equation}
where $V$ is the relative speed at the quarter-chord location, and
\begin{equation}
C_n(y)=2\pi\left(\alpha^\prime+\alpha_0+\bar{\theta}\right)
\tag{2}\label{eq:2}
\end{equation}
with $\alpha_0$ the section zero-lift angle and $\bar{\theta}$ the section mean pitch angle. The mean pitch is
\begin{equation}
\bar{\theta}=\bar{\theta}_a+\bar{\theta}_w
\tag{3}\label{eq:3}
\end{equation}
where $\bar{\theta}_a$ is the flapping-axis angle relative to the freestream, and $\bar{\theta}_w$ is the mean chord
angle relative to the flapping axis.

\subsection{Finite-span unsteady correction}
The effective unsteady term $\alpha^\prime$ is written as
\begin{equation}
\alpha^\prime=\left[\frac{AR\,C^\prime(k)}{2+AR}\right]\alpha-\frac{w_0}{U}
\tag{4}\label{eq:4}
\end{equation}
where $\alpha$ is the relative angle of attack at the $3/4$-chord location due to the motion:
\begin{equation}
\alpha=\frac{\dot{h}\cos(\theta-\bar{\theta}_a)+\frac{3}{4}c\dot{\theta}+U(\theta-\bar{\theta})}{U}
\tag{5}\label{eq:5}
\end{equation}

In the strip-theory interpretation adopted here, each chordwise strip is assumed to behave as if it were part of an
equivalent elliptical-planform wing of the same aspect ratio, executing a simple harmonic motion consistent with the
strip kinematics. For a finite-span wing, the unsteady normal-force coefficient increment may be expressed as
\begin{equation}
\delta C_n=2\pi\,C(k)_{\text{Jones}}\,\alpha
\tag{6}\label{eq:6}
\end{equation}
where the reduced frequency is
\begin{equation}
k=\frac{c\omega}{2U}
\tag{7}\label{eq:7}
\end{equation}
and a convenient representation is
\begin{equation}
C(k)_{\text{Jones}}=\frac{AR\,C^\prime(k)}{2+AR}
\tag{8}\label{eq:8}
\end{equation}
with the complex function
\begin{equation}
C^\prime(k)=F^\prime(k)+iG^\prime(k).
\tag{9}\label{eq:9}
\end{equation}

An approximation for the complex terms is
\begin{align}
F^\prime(k)&=1-\frac{C_1k^2}{k^2+C_2^2}\notag\\
G^\prime(k)&=-\frac{C_1C_2k}{k^2+C_2^2}\notag\\
C_1&=\frac{0.5\,AR}{2.32+AR}\notag\\
C_2&=0.181+\frac{0.772}{AR}.
\tag{9}\label{eq:9_block}
\end{align}

Assuming harmonic motion,
\begin{equation}
\alpha=Ae^{i\omega t}
\tag{10}\label{eq:10}
\end{equation}
one obtains
\begin{equation}
\alpha^\prime=\frac{AR}{2+AR}\left[F^\prime(k)\alpha+\frac{c}{2U}\frac{G^\prime(k)}{k}\dot{\alpha}\right]-\frac{w_0}{U}.
\tag{11}\label{eq:11}
\end{equation}

A compact approximation for downwash is
\begin{equation}
\frac{w_0}{U}=\frac{2(\alpha_0+\bar{\theta})}{2+AR}.
\tag{12}\label{eq:12}
\end{equation}
If $(\alpha_0+\bar{\theta})$ varies significantly along the span, a more accurate evaluation of $w_0/U$
may be required (e.g., a lifting-line type correction for twisted wings), beyond the compact estimate in \eqref{eq:12}.

\subsection{Relative speed and apparent-mass normal force}
The relative speed at the quarter-chord is
\begin{equation}
V=\left\{\left[U\cos\theta-\dot{h}\sin(\theta-\bar{\theta}_a)\right]^2+
\left[U(\alpha^\prime+\bar{\theta})-\frac{1}{2}c\dot{\theta}\right]^2\right\}^{1/2}
\tag{13}\label{eq:13}
\end{equation}

The apparent-mass contribution (acting at midchord) is
\begin{equation}
dN_a=\frac{\rho\pi c^2}{4}\,\dot{v}_2\,dy
\tag{14}\label{eq:14}
\end{equation}
where the midchord normal-velocity rate is
\begin{equation}
\dot{v}_2=U\dot{\alpha}-\frac{1}{4}c\ddot{\theta}.
\tag{15}\label{eq:15}
\end{equation}
Thus the total attached-flow normal force is
\begin{equation}
dN=dN_c+dN_a.
\tag{16}\label{eq:16}
\end{equation}

\subsection{Chordwise forces: camber, leading-edge suction, and skin friction}
The chordwise force due to camber is
\begin{equation}
dD_{\text{camber}}=-2\pi\alpha_0(\alpha^\prime+\bar{\theta})\left(\frac{\rho U V}{2}\right)c\,dy
\tag{17}\label{eq:17}
\end{equation}

The leading-edge suction contribution (with efficiency $\eta_s$) is
\begin{equation}
dT_s=\eta_s\,2\pi\left(\alpha^\prime+\bar{\theta}-\frac{1}{4}\frac{c\dot{\theta}}{U}\right)^2
\left(\frac{\rho U V}{2}\right)c\,dy
\tag{18}\label{eq:18}
\end{equation}

Viscosity introduces chordwise skin-friction drag
\begin{equation}
dD_f=\left(C_d\right)_f\left(\frac{\rho V_x^2}{2}\right)c\,dy
\tag{19}\label{eq:19}
\end{equation}
where the tangential speed is approximated by
\begin{equation}
V_x=U\cos\theta-\dot{h}\sin(\theta-\bar{\theta}_a).
\tag{20}\label{eq:20}
\end{equation}
Hence the total chordwise force is
\begin{equation}
dF_x=dT_s-dD_{\text{camber}}-dD_f.
\tag{21}\label{eq:21}
\end{equation}

\subsection{Dynamic stall criterion and separated-flow model}
A localized dynamic-stall angle is modeled as
\begin{equation}
\alpha_{\text{stall}}=\left(\alpha_{\text{stall}}\right)_{\text{static}}+\xi\left[\frac{c\dot{\alpha}}{2U}\right]^{1/2}
\tag{22}\label{eq:22}
\end{equation}
and applied at the leading edge. The attached-flow condition is
\begin{equation}
\left(\alpha_{\text{stall}}\right)_{\min}\le
\left[\alpha^\prime+\bar{\theta}-\frac{3}{4}\left(\frac{c\dot{\theta}}{U}\right)\right]\le
\left(\alpha_{\text{stall}}\right)_{\max}.
\tag{23}\label{eq:23}
\end{equation}

When the attached-flow range is exceeded, totally separated flow is assumed to occur abruptly,
in which case all chordwise forces are neglected:
\begin{equation}
dD_{\text{camber}},\,dT_s,\,dD_f=0.
\tag{24}\label{eq:24}
\end{equation}
The normal force is then
\begin{equation}
dN=\left(dN_c\right)_{\text{sep}}+\left(dN_a\right)_{\text{sep}}.
\tag{25}\label{eq:25}
\end{equation}
The separated-flow circulatory normal force (crossflow drag) is
\begin{equation}
\left(dN_c\right)_{\text{sep}}=\left(C_d\right)_{\text{cf}}\left(\frac{\rho\hat{V}V_n}{2}\right)c\,dy
\tag{26}\label{eq:26}
\end{equation}
where
\begin{equation}
\hat{V}=\left(V_x^2+V_n^2\right)^{1/2}
\tag{27}\label{eq:27}
\end{equation}
and the midchord normal velocity component is
\begin{equation}
V_n=\dot{h}\cos(\theta-\bar{\theta}_a)+\frac{1}{2}c\dot{\theta}+U\sin\theta.
\tag{28}\label{eq:28}
\end{equation}
Note that $\dot{v}_2$ in \eqref{eq:15} corresponds to a linearized time derivative associated with $V_n$.
The apparent-mass term in separation is taken as
\begin{equation}
\left(dN_a\right)_{\text{sep}}=\frac{1}{2}dN_a.
\tag{29}\label{eq:29}
\end{equation}

\subsection{Instantaneous lift and thrust; span integration; cycle averaging}
The strip lift and thrust are
\begin{equation}
dL=dN\cos\theta+dF_x\sin\theta
\tag{30}\label{eq:30}
\end{equation}
\begin{equation}
dT=dF_x\cos\theta-dN\sin\theta.
\tag{31}\label{eq:31}
\end{equation}

Integrating along the span gives the whole-wing instantaneous lift and thrust:
\begin{equation}
L(t)=2\int_{0}^{b/2}\cos\gamma\,dL,\qquad
T(t)=2\int_{0}^{b/2}dT
\tag{32}\label{eq:32}
\end{equation}
where $\gamma(t)$ is the instantaneous dihedral angle of the section.

It is convenient to integrate over cycle angle $\phi$ instead of time:
\begin{equation}
\phi=\omega t
\tag{33}\label{eq:33}
\end{equation}
so that the cycle-averaged lift and thrust are
\begin{equation}
\bar{L}=\frac{1}{2\pi}\int_{0}^{2\pi}L(\phi)\,d\phi,\qquad
\bar{T}=\frac{1}{2\pi}\int_{0}^{2\pi}T(\phi)\,d\phi.
\tag{34}\label{eq:34}
\end{equation}

\subsection{Input power, output power, and efficiency}
For attached flow, the instantaneous power required to move the section against aerodynamic loads is
\begin{align}
dP_{\text{in}}&=
dF_x\,\dot{h}\sin(\theta-\bar{\theta}_a)
+dN\left[\dot{h}\cos(\theta-\bar{\theta}_a)+\frac{1}{4}c\dot{\theta}\right]
+dN_a\left[\frac{1}{4}c\dot{\theta}\right]
-dM_{ac}\dot{\theta}-dM_a\dot{\theta}.
\tag{35}\label{eq:35}
\end{align}
Here $dM_{ac}$ is the pitching moment about the aerodynamic centre and $dM_a$ includes apparent-camber and apparent-inertia moments:
\begin{equation}
dM_a=-\left[\frac{1}{16}\rho\pi c^3\dot{\theta}U+\frac{1}{128}\rho\pi c^4\ddot{\theta}\right]dy.
\tag{36}\label{eq:36}
\end{equation}

For separated flow, the input power expression becomes
\begin{equation}
dP_{\text{in}}=dN_{\text{sep}}\left[\dot{h}\cos(\theta-\bar{\theta}_a)+\frac{1}{2}c\dot{\theta}\right].
\tag{37}\label{eq:37}
\end{equation}

The instantaneous aerodynamic power absorbed by the whole wing is
\begin{equation}
P_{\text{in}}(t)=2\int_{0}^{b/2}dP_{\text{in}}
\tag{38}\label{eq:38}
\end{equation}
and the average input power over the cycle is
\begin{equation}
\bar{P}_{\text{in}}=\frac{1}{2\pi}\int_{0}^{2\pi}P_{\text{in}}(\phi)\,d\phi.
\tag{39}\label{eq:39}
\end{equation}

The average output power is
\begin{equation}
\bar{P}_{\text{out}}=\bar{T}U
\tag{40}\label{eq:40}
\end{equation}
and the average propulsive efficiency is
\begin{equation}
\bar{\eta}=\frac{\bar{P}_{\text{out}}}{\bar{P}_{\text{in}}}.
\tag{41}\label{eq:41}
\end{equation}

% =========================================================
\section{Numerical Example (parameter and discretization statements)}
This section records the example parameter choices and discretizations in a compact, reproducible form.

\subsection{Example parameter set}
In the reported example setup, negative-$\alpha^\prime$ stalling is assumed not to occur,
so $(\alpha_{\text{stall}})_{\min}$ is not specified. In addition, the dynamic-stall criterion
was not effectively used in the calculations at the time, so the parameter $\xi$ in \eqref{eq:22}
is taken as zero for the example.

A representative parameter set is
\begin{align}
\alpha_0&=0.5^\circ\notag\\
\eta_s&=0.98\notag\\
C_{mac}&=0.025\notag\\
\left(\alpha_{\text{stall}}\right)_{\max}&=13^\circ.
\tag{42}\label{eq:42}
\end{align}

For the separated-flow normal-force model, a high-aspect-ratio flat-plate value is used:
$(C_d)_{\text{cf}}=1.98$ (example setting). The assumed kinematics also fix the phasing between
plunging and pitching at $-90^\circ$.

\subsection{Skin-friction drag coefficient (example)}
A friction-drag coefficient for a fully turbulent boundary layer is taken as
\begin{equation}
\left(C_d\right)_f=\frac{0.89}{\left[\log(R_n)\right]^{2.58}}
\tag{43}\label{eq:43}
\end{equation}
where $R_n$ is based on the local chord.

\subsection{Root-flapping kinematics used in the example}
With root-flapping kinematics and no spanwise bending, the plunging motion is
\begin{equation}
h=-(\Gamma y)\cos\phi
\tag{44}\label{eq:44}
\end{equation}
where $\Gamma$ is the maximum flapping-angle magnitude (e.g., $\Gamma=0.3491~\text{rad}$).

A linear spanwise dynamic twist distribution is assumed:
\begin{equation}
\delta\theta=-(\beta_0 y)\sin\phi
\tag{45}\label{eq:45}
\end{equation}
where $\beta_0$ is a chosen dynamic-twist parameter.

\subsection{Discrete integration (cycle and span)}
A discrete approximation of the cycle integral (example shown for $\bar{L}$) is
\begin{equation}
\bar{L}\cong\frac{1}{m}\sum_{j=1}^{m}L_j
\tag{46}\label{eq:46}
\end{equation}
and, using the span integral structure of \eqref{eq:32}, one may approximate
\begin{equation}
L_j\cong 2\sum_{i=1}^{n}\cos\gamma_{ij}\,\Delta L_{ij}.
\tag{47}\label{eq:47}
\end{equation}

% =========================================================
\section{Notes for implementation}
To reproduce results numerically:
\begin{itemize}
\item Discretize span ($i=1\ldots n$) and cycle phase ($j=1\ldots m$).
\item For each $(i,j)$, compute $c(y_i)$, kinematics $(\theta,\dot{\theta},\ddot{\theta},h,\dot{h})$ and derived
quantities $(\alpha,\dot{\alpha},\alpha^\prime,V,V_x,V_n)$ using \eqref{eq:4}--\eqref{eq:28}.
\item Apply attached/separated switching via \eqref{eq:23}; compute forces with \eqref{eq:16}--\eqref{eq:31}.
\item Compute power via \eqref{eq:35} or \eqref{eq:37}, then integrate with \eqref{eq:38}--\eqref{eq:41} and
(optional) discrete rules \eqref{eq:46}--\eqref{eq:47}.
\item If the time history is not uniform in phase (e.g., downstroke faster than upstroke), perform time-weighted
averaging using the correct $dt$ rather than uniform $\phi$ sampling.
\end{itemize}

% =========================================================
\section{Conclusions (paraphrased)}
A finite-span, unsteady strip-theory framework can be assembled to estimate mean lift, thrust, and power of
harmonically flapping wings while accounting for leading-edge suction efficiency, friction drag, and localized
stall/separation effects. A central assumption is that the semispan remains constant throughout the motion;
the example application indicates that sustained flight may depend strongly on an appropriate dynamic twist distribution.

\section*{Acknowledgements (paraphrased)}
This rewrite omits the original reference list by request.

\end{document}
